% 02_h05_volumetric_printing
% Two volumetric photopolymerisation principles, top views.
% Left: xolography (xolo; Regehly et al., Nature 588, 620-624, 2020): a UV light sheet switches a
%   dual-colour photoinitiator into a latent state; a projected visible image cures the resin only
%   where it meets the sheet; the sheet sweeps through the resin, no layers, no supports.
% Right: tomographic printing (Readily3D; Kelly et al., Science 363, 1075, 2019;
%   Loterie et al., Nat. Commun. 11, 852, 2020): the vial rotates, the projector sends the matching
%   image at every angle, the light dose adds up and the resin cures where it passes a threshold.
\documentclass[tikz,border=8pt]{standalone}
\usepackage{xfab-tikz}
\begin{document}
\begin{tikzpicture}[x=1cm,y=1cm]
% ---------------------------------------------------------------- xolography
\node[font=\bfseries] at (2,5.15) {xolography};
\node[font=\footnotesize, text=inksoft] at (2,4.72) {two lights cross};
% cuvette, top view
\fill[cIFeed!30] (0,0) rectangle (4,4);
\draw[cToolD, line width=1.2pt] (0,0) rectangle (4,4);
% printed part (far side, already swept)
\fill[ipart] (1.0,2.35) rectangle (3.0,3.35);
\fill[cIFeed!30] (1.6,2.35) rectangle (2.4,2.95);
\draw[cIPartD, line width=0.7pt] (1.6,2.35) -- (1.6,2.95) -- (2.4,2.95) -- (2.4,2.35);
% projected visible image: cone from the projector below
\fill[cLight!12] (1.55,-1.15) -- (2.45,-1.15) -- (3.0,2.2) -- (1.0,2.2) -- cycle;
\draw[cLight!55, line width=0.6pt, dash pattern=on 3pt off 2pt] (1.55,-1.15) -- (1.0,2.2);
\draw[cLight!55, line width=0.6pt, dash pattern=on 3pt off 2pt] (2.45,-1.15) -- (3.0,2.2);
\fill[tool] (1.3,-1.75) rectangle (2.7,-1.15);
\node[font=\scriptsize] at (2,-1.45) {projector};
% UV light sheet from the left
\fill[tool] (-1.9,1.98) rectangle (-0.9,2.42);
\node[font=\scriptsize, align=center] at (-1.4,2.75) {UV line\\[-1pt]laser};
\draw[lighthalo] (-0.9,2.2) -- (4,2.2);
\draw[light] (-0.9,2.2) -- (4,2.2);
% curing slice where both lights meet
\draw[cIPartD, line width=2.4pt] (1.0,2.2) -- (1.6,2.2);
\draw[cIPartD, line width=2.4pt] (2.4,2.2) -- (3.0,2.2);
% sweep of the sheet
\draw[move] (4.35,2.55) -- (4.35,1.55);
\node[lab, anchor=west] at (4.45,2.05) {sheet\\sweeps};
% labels
\node[lab, anchor=west] (a) at (4.45,3.45) {printed part,\\held by the resin};
\draw[lead] (a.west) -- (2.95,3.2);
\node[lab, anchor=west] (b) at (4.45,0.75) {resin cures only\\where image\\and sheet meet};
\draw[lead] (b.west) -- (2.75,2.2);
\node[lab, anchor=north] at (2,-1.9) {$\lambda_1$ UV sheet switches the initiator on,\\$\lambda_2$ visible image cures the slice};

% ---------------------------------------------------------------- tomographic
\begin{scope}[shift={(10.2,0)}]
\node[font=\bfseries] at (2,5.15) {tomographic};
\node[font=\footnotesize, text=inksoft] at (2,4.72) {one light, all angles};
% vial, top view
\fill[cIFeed!30] (2,2) circle (1.6);
\draw[cToolD, line width=1.2pt] (2,2) circle (1.6);
% part in the centre
\fill[ipart] (1.45,1.4) rectangle (2.55,2.6);
\fill[cIFeed!30] (1.85,1.4) rectangle (2.55,2.15);
\draw[cIPartD, line width=0.7pt] (1.85,1.4) -- (1.85,2.15) -- (2.55,2.15);
% projected pattern from the left: parallel rays with gaps
\fill[tool] (-1.9,1.25) rectangle (-1.2,2.75);
\node[font=\scriptsize, rotate=90] at (-1.55,2.0) {projector};
\foreach \y/\x in {1.45/3.55, 1.65/3.6, 1.85/3.6, 2.35/3.6, 2.55/3.55}{
  \draw[cLight!70, line width=1.1pt] (-1.2,\y) -- (\x,\y);}
% rotation
\draw[move, line width=1pt] (2,2) ++(40:1.95) arc (40:140:1.95);
\node[lab, anchor=south] at (2,3.95) {vial rotates $360^\circ$};
% labels
\node[lab, anchor=west] (c) at (4.0,0.55) {dose adds up:\\cured where it\\passes a threshold};
\draw[lead] (c.west) -- (1.65,1.75);
\node[lab, anchor=north] at (2,-1.9) {the projector shows the matching\\image at every angle};
\end{scope}
\end{tikzpicture}
\end{document}
